% ===== Paquetes =====
\usepackage{xcolor}
\usepackage{titlesec}
\usepackage{fancyhdr}
\usepackage{lastpage}
\usepackage[most]{tcolorbox}
\usepackage{booktabs}
\usepackage{caption}
\usepackage{float}          % opción [H]: fija la figura donde se declara
\usepackage{enumitem}
% Pandoc carga parskip ANTES de este preámbulo y parskip redefine \@starttoc;
% eso hacía que tocloft se inhibiera ("tocloft bailing out") y el TOC perdiera
% el formato institucional. Se restaura la definición del kernel y luego se carga tocloft.
% (la definición debe coincidir token a token con la del kernel: \def, no
%  \newcommand, porque tocloft la compara con \ifx)
\makeatletter
\def\@starttoc#1{%
 \begingroup
   \makeatletter
   \@input{\jobname.#1}%
   \if@filesw
     \expandafter\newwrite\csname tf@#1\endcsname
     \immediate\openout \csname tf@#1\endcsname \jobname.#1\relax
   \fi
  \@nobreakfalse
\endgroup}
\makeatother
\usepackage{tocloft}
\usepackage{parskip}          % espacio entre párrafos, sin sangría
\usepackage{tikz}
\usetikzlibrary{calc, positioning, arrows.meta, fit, backgrounds}

% ===== Colores institucionales — Universidad Externado de Colombia =====
% Tomados de la hoja de estilos vigente del sitio oficial (uexternado.edu.co,
% tema WordPress "externado_theme" + pattern-lab-externado-prod), no de un
% manual de marca publicado (no se localizó uno público con códigos Pantone).
\definecolor{extVerde}{HTML}{4D7B19}      % verde institucional primario (bordes, acentos, enlaces)
\definecolor{extVerdeClaro}{HTML}{78B92F} % verde secundario (enlaces sobre fondo oscuro, resaltes)
\definecolor{extVerdeOscuro}{HTML}{00522C} % verde profundo (bandas, contraste fuerte)
\definecolor{extDorado}{HTML}{E7A900}     % dorado/mostaza (del escudo/bandera institucional)
\definecolor{extNegro}{HTML}{1A1A1A}
\definecolor{extGris}{HTML}{58595B}       % gris neutro para texto secundario (no tomado del sitio;
                                           % elección propia, análoga al rol de ipnGris)

% Colores de apoyo para diagramas (armonizados con la paleta institucional)
\definecolor{diagAzul}{HTML}{2F4B7C}
\definecolor{diagDorado}{HTML}{B5810F}
\definecolor{diagVerdeOscuro}{HTML}{00522C}

% ===== Estilos de diagrama de flujo (genéricos, listos para usar) =====
\tikzset{
  paso/.style={
    draw, rounded corners=3pt, line width=0.7pt,
    text width=6.1cm, align=center, inner sep=6pt, font=\sffamily\small,
  },
  % (el título del paso se compone con \pasotitulo, definido abajo)
  neutro/.style={paso, draw=extGris,       fill=extGris!7,       text=extNegro},
  principal/.style={paso, draw=extVerde,   fill=extVerde!7,      text=extVerde},
  secundario/.style={paso, draw=diagAzul,  fill=diagAzul!7,      text=diagAzul},
  destacado/.style={paso, draw=diagDorado, fill=diagDorado!10,   text=diagDorado},
  final/.style={paso, draw=diagVerdeOscuro, fill=diagVerdeOscuro!8, text=diagVerdeOscuro},
  etapa/.style={font=\sffamily\footnotesize\bfseries, text=extNegro, anchor=west},
  flujo/.style={-{Stealth[length=2.4mm]}, draw=extGris, line width=0.7pt},
  clave/.style={draw, rounded corners=2pt, line width=0.6pt,
                font=\sffamily\scriptsize, inner sep=3pt, text width=2.5cm, align=center},
}
% Título en negrita dentro de una caja del diagrama
\newcommand{\pasotitulo}[1]{{\bfseries #1}}

% Color de texto base
\color{extNegro}

% ===== Títulos (Lato vía \sffamily + verde) =====
% "action titles": van en \section como oraciones completas
\titleformat{\section}
  {\sffamily\bfseries\large\color{extVerde}}{\thesection}{0.6em}{}
\titleformat{\subsection}
  {\sffamily\bfseries\normalsize\color{extNegro}}{\thesubsection}{0.6em}{}
\titleformat{\subsubsection}
  {\sffamily\bfseries\small\color{extGris}}{}{0em}{}
\titlespacing*{\section}{0pt}{14pt}{5pt}
\titlespacing*{\subsection}{0pt}{10pt}{3pt}

% ===== Pie de página con numeración =====
\pagestyle{fancy}
\fancyhf{}
\renewcommand{\headrulewidth}{0pt}
\renewcommand{\footrulewidth}{0.4pt}
\renewcommand{\footrule}{\hbox to\headwidth{\color{extGris!40}\leaders\hrule height \footrulewidth\hfill}}
\fancyfoot[L]{\sffamily\footnotesize\color{extGris}%
  Universidad Externado de Colombia · Informe de Deep Learning}
\fancyfoot[R]{\sffamily\footnotesize\color{extGris}%
  Pág. \thepage\ de \pageref{LastPage}}



% (tocloft se carga arriba, antes de parskip)
\renewcommand{\cftsecfont}{\sffamily\bfseries\color{extNegro}}
\renewcommand{\cftsecpagefont}{\sffamily\color{extGris}}
\renewcommand{\cftsubsecfont}{\sffamily\color{extGris}}
\renewcommand{\cftsubsecpagefont}{\sffamily\color{extGris}}
\renewcommand{\cfttoctitlefont}{\sffamily\bfseries\LARGE\color{extVerde}}
\setlength{\cftbeforesecskip}{6pt}
% Etiquetas anchas tipo "Anexo A" / "Anexo C.1": la caja del número por defecto
% (~2.3em) es muy angosta y la etiqueta se sale y se solapa con el título. Se ensancha:
\setlength{\cftsecnumwidth}{4.7em}      % número de sección: cabe "Anexo A"
\setlength{\cftsubsecindent}{4.7em}     % subsección alineada bajo el título de sección
\setlength{\cftsubsecnumwidth}{5.6em}   % número de subsección: cabe "Anexo C.1"

% ===== Pies de figura/tabla =====
\captionsetup{font=small, labelfont={sf,bf,color=extVerde}, textfont={sf}}
% babel-spanish rotula las tablas como "Cuadro"; se prefiere "Tabla",
% que además es como se las nombra habitualmente en el cuerpo del texto.
\addto\captionsspanish{\renewcommand{\tablename}{Tabla}}

% ===== Caja "pensamiento gobernante" / callouts =====
\newtcolorbox{govbox}{
  enhanced, breakable, boxrule=0pt, leftrule=3pt,
  colback=extVerde!6, colframe=extVerde,
  arc=1pt, left=10pt, right=10pt, top=8pt, bottom=8pt,
  fontupper=\normalsize
}

% ===== Metadatos del PDF =====
\usepackage{hyperref}
\hypersetup{
  pdftitle={Informe de Deep Learning — Universidad Externado de Colombia},
  pdfauthor={Yeiver Estevan Sanabria Sánchez},
  colorlinks=true, linkcolor=extVerde, urlcolor=extVerde, citecolor=extVerde
}
